%%%%%%%%%%%%%%%%
%%% deut 18 %%%
%%%%%%%%%%%%%%%%

\Chapter{18}

\Verse{1}
{
	\hDtn{לֹא יִהְיֶה לַכֹּהֲנִים הַלְוִיִּם כּל שֵׁבֶט לֵוִי חֵלֶק וְנַחֲלָה עִם יִשְׂרָאֵל אִשֵּׁי יְהֹוָה וְנַחֲלָתוֹ יֹאכֵלוּן׃}
}
{
	\eDtn{
		“The Levite priests, all the tribe of Levi, shall not
		have a portion and legacy with Israel. They shall eat
		YHWH’s offerings by fire and His legacy.
	}
}
{
}

\Verse{2}
{
	\hDtn{וְנַחֲלָה לֹא יִהְיֶה לּוֹ בְּקֶרֶב אֶחָיו יְהֹוָה הוּא נַחֲלָתוֹ כַּאֲשֶׁר דִּבֶּר לוֹ׃}
}
{
	\eDtn{
		But he shall not have a legacy among his brothers. YHWH:
		He is his legacy, as He spoke to him.
	}
}
{
}

\Verse{3}
{
	\hDtn{וְזֶה יִהְיֶה מִשְׁפַּט הַכֹּהֲנִים מֵאֵת הָעָם מֵאֵת זֹבְחֵי הַזֶּבַח אִם שׁוֹר אִם שֶׂה וְנָתַן לַכֹּהֵן הַזְּרֹעַ וְהַלְּחָיַיִם וְהַקֵּבָה׃}
}
{
	\eDtn{
		And this shall be the rule for the priests from the
		people, from those who make a sacrifice, whether an ox or
		a sheep: he shall give the shoulder, the cheeks, and the
		stomach.
	}
}
{
}

\Verse{4}
{
	\hDtn{רֵאשִׁית דְּגָנְךָ תִּירֹשְׁךָ וְיִצְהָרֶךָ וְרֵאשִׁית גֵּז צֹאנְךָ תִּתֶּן לוֹ׃}
}
{
	\eDtn{
		You shall give him the first of your grain, your wine, and
		your oil, and the first shearing of your sheep.
	}
}
{
}

\Verse{5}
{
	\hDtn{כִּי בוֹ בָּחַר יְהֹוָה אֱלֹהֶיךָ מִכּל שְׁבָטֶיךָ לַעֲמֹד לְשָׁרֵת בְּשֵׁם יְהֹוָה הוּא וּבָנָיו כּל הַיָּמִים׃}
}
{
	\eDtn{
		Because YHWH, your God, has chosen him from all your
		tribes to stand to serve in YHWH’s name, he and his sons,
		for all time.
	}
}
{
}

\Verse{6}
{
	\hDtn{וְכִי יָבֹא הַלֵּוִי מֵאַחַד שְׁעָרֶיךָ מִכּל יִשְׂרָאֵל אֲשֶׁר הוּא גָּר שָׁם וּבָא בְּכל אַוַּת נַפְשׁוֹ אֶל הַמָּקוֹם אֲשֶׁר יִבְחַר יְהֹוָה׃}
}
{
	\eDtn{
		“And when a Levite will come from one of your gates, from
		all of Israel, where he lives, then he shall come as much
		as his soul desires to the place that YHWH will choose
	}
}
{
}

\Verse{7}
{
	\hDtn{וְשֵׁרֵת בְּשֵׁם יְהֹוָה אֱלֹהָיו כְּכל אֶחָיו הַלְוִיִּם הָעֹמְדִים שָׁם לִפְנֵי יְהֹוָה׃}
}
{
	\eDtn{
		and serve in the name of YHWH, his God, like all of his
		Levite brothers who are standing there in front of YHWH.
	}
}
{
}

\Verse{8}
{
	\hDtn{חֵלֶק כְּחֵלֶק יֹאכֵלוּ לְבַד מִמְכָּרָיו עַל הָאָבוֹת׃}
}
{
	\eDtn{
		They shall eat, portion for portion, aside from his sales
		of patrimony.
	}
}
{
%	\footnote{
%		portion for portion. Any Levite who comes to serve at the central place
%		of worship is supposed to receive the same share of the priestly food as
%		every other Levite. And this share is not affected by whether he has any
%		other income from family property, i.e., “from his sales of patrimony.”
%	}
}

\Verse{9}
{
	\hDtn{כִּי אַתָּה בָּא אֶל הָאָרֶץ אֲשֶׁר יְהֹוָה אֱלֹהֶיךָ נֹתֵן לָךְ לֹא תִלְמַד לַעֲשׂוֹת כְּתוֹעֲבֹת הַגּוֹיִם הָהֵם׃}
}
{
	\eDtn{
		“When you come to the land that YHWH, your God, is giving
		you, you shall not learn to do like the offensive things
		of those nations.
	}
}
{
}

\Verse{10}
{
	\hDtn{לֹא יִמָּצֵא בְךָ מַעֲבִיר בְּנוֹ וּבִתּוֹ בָּאֵשׁ קֹסֵם קְסָמִים מְעוֹנֵן וּמְנַחֵשׁ וּמְכַשֵּׁף׃}
}
{
	\eDtn{
		There shall not be found among you someone who passes his
		son or his daughter through fire, one who practices
		enchantment, a soothsayer or a diviner or a sorcerer
	}
}
{
}

\Verse{11}
{
	\hDtn{וְחֹבֵר חָבֶר וְשֹׁאֵל אוֹב וְיִדְּעֹנִי וְדֹרֵשׁ אֶל הַמֵּתִים׃}
}
{
	\eDtn{
		or one who casts spells or who asks of a ghost or of a
		spirit of an acquaintance or inquires of the dead,
	}
}
{
}

\Verse{12}
{
	\hDtn{כִּי תוֹעֲבַת יְהֹוָה כּל עֹשֵׂה אֵלֶּה וּבִגְלַל הַתּוֹעֵבֹת הָאֵלֶּה יְהֹוָה אֱלֹהֶיךָ מוֹרִישׁ אוֹתָם מִפָּנֶיךָ׃}
}
{
	\eDtn{
		because everyone who does these is an offensive thing of
		YHWH, and because of these offensive things YHWH, your
		God, is dispossessing them from in front of you.
	}
}
{
}

\Verse{13}
{
	\hDtn{תָּמִים תִּהְיֶה עִם יְהֹוָה אֱלֹהֶיךָ׃}
}
{
	\eDtn{You shall be unblemished with YHWH, your God,}
}
{
%	\footnote{
%		unblemished. This word is the mark of the recipients of the first two
%		covenants with God. It is the word that is used to describe the merit of
%		Noah; and God tells Abraham to be unblemished as well in the opening
%		words of the covenant (Gen 6:9; 17:1). (The only other person in the
%		Tanak to be described by these words is Job.) Now all the people of
%		Israel are recipients of the third covenant, and so all of them are
%		instructed to be like Noah and Abraham. They may receive the covenant
%		initially on the merit of Abraham, but they must maintain it on their
%		own integrity.
%	}
}

\Verse{14}
{
	\hDtn{כִּי הַגּוֹיִם הָאֵלֶּה אֲשֶׁר אַתָּה יוֹרֵשׁ אוֹתָם אֶל מְעֹנְנִים וְאֶל קֹסְמִים יִשְׁמָעוּ וְאַתָּה לֹא כֵן נָתַן לְךָ יְהֹוָה אֱלֹהֶיךָ׃}
}
{
	\eDtn{
		because these nations whom you are dispossessing listen to
		soothsayers and enchanters, but you: YHWH, your God, has
		not permitted such for you.
	}
}
{
}

\Verse{15}
{
	\hDtn{נָבִיא מִקִּרְבְּךָ מֵאַחֶיךָ כָּמֹנִי יָקִים לְךָ יְהֹוָה אֱלֹהֶיךָ אֵלָיו תִּשְׁמָעוּן׃}
}
{
	\eDtn{
		YHWH, your God, will raise up for you a prophet from among
		you, from your brothers, like me. You shall listen to him—
	}
}
{
%	\footnote{
%		God will raise up a prophet like me. But the end of the Torah says: “A
%		prophet did not rise again in Israel like Moses” (Deut 34:10)! Still,
%		this need not be a contradiction. The context of the passage here in
%		Deuteronomy 18 is that Israel should not listen to soothsayers and
%		enchanters, but only to a prophet from God, a “prophet like me.” Moses
%		reminds the people that this is what they asked for at Sinai: that they
%		not hear God’s voice directly anymore, but that they hear from God only
%		through prophets. The passage at the end of the Torah visibly means that
%		no other prophet was as great as Moses. It is simply a linguistic
%		matter: the range of the expression “to be like someone” is wide enough
%		to have these two meanings (and several more).
%	}
}

\Verse{16}
{
	\hDtn{כְּכֹל אֲשֶׁר שָׁאַלְתָּ מֵעִם יְהֹוָה אֱלֹהֶיךָ בְּחֹרֵב בְּיוֹם הַקָּהָל לֵאמֹר לֹא אֹסֵף לִשְׁמֹעַ אֶת קוֹל יְהֹוָה אֱלֹהָי וְאֶת הָאֵשׁ הַגְּדֹלָה הַזֹּאת לֹא אֶרְאֶה עוֹד וְלֹא אָמוּת׃}
}
{
	\eDtn{
		in accordance with everything that you asked from YHWH,
		your God, at Horeb in the day of the assembly, saying,
		‘Let me not continue to hear the voice of YHWH, my God,
		and let me not see this big fire anymore, so I won’t die!’
	}
}
{
}

\Verse{17}
{
	\hDtn{וַיֹּאמֶר יְהֹוָה אֵלָי הֵיטִיבוּ אֲשֶׁר דִּבֵּרוּ׃}
}
{
	\eDtn{
		“And YHWH said to me, ‘They’ve been good in what they’ve
		spoken.
	}
}
{
}

\Verse{18}
{
	\hDtn{נָבִיא אָקִים לָהֶם מִקֶּרֶב אֲחֵיהֶם כָּמוֹךָ וְנָתַתִּי דְבָרַי בְּפִיו וְדִבֶּר אֲלֵיהֶם אֵת כּל אֲשֶׁר אֲצַוֶּנּוּ׃}
}
{
	\eDtn{
		I’ll raise up a prophet for them from among their
		brothers, like you, and I’ll put my words in his mouth,
		and he’ll speak to them everything that I’ll command him.
	}
}
{
%	\footnote{
%		God will raise up a prophet like me. But the end of the Torah says: “A
%		prophet did not rise again in Israel like Moses” (Deut 34:10)! Still,
%		this need not be a contradiction. The context of the passage here in
%		Deuteronomy 18 is that Israel should not listen to soothsayers and
%		enchanters, but only to a prophet from God, a “prophet like me.” Moses
%		reminds the people that this is what they asked for at Sinai: that they
%		not hear God’s voice directly anymore, but that they hear from God only
%		through prophets. The passage at the end of the Torah visibly means that
%		no other prophet was as great as Moses. It is simply a linguistic
%		matter: the range of the expression “to be like someone” is wide enough
%		to have these two meanings (and several more).
%	}
}

\Verse{19}
{
	\hDtn{וְהָיָה הָאִישׁ אֲשֶׁר לֹא יִשְׁמַע אֶל דְּבָרַי אֲשֶׁר יְדַבֵּר בִּשְׁמִי אָנֹכִי אֶדְרֹשׁ מֵעִמּוֹ׃}
}
{
	\eDtn{
		And it will be: the man who won’t listen to my words that
		he’ll speak in my name, I shall require it from him!
	}
}
{
}

\Verse{20}
{
	\hDtn{אַךְ הַנָּבִיא אֲשֶׁר יָזִיד לְדַבֵּר דָּבָר בִּשְׁמִי אֵת אֲשֶׁר לֹא צִוִּיתִיו לְדַבֵּר וַאֲשֶׁר יְדַבֵּר בְּשֵׁם אֱלֹהִים אֲחֵרִים וּמֵת הַנָּבִיא הַהוּא׃}
}
{
	\eDtn{
		Just: the prophet who will presume to speak a thing in my
		name that I didn’t command him to speak, and who will
		speak in the name of other gods—that prophet shall die.’
	}
}
{
}

\Verse{21}
{
	\hDtn{וְכִי תֹאמַר בִּלְבָבֶךָ אֵיכָה נֵדַע אֶת הַדָּבָר אֲשֶׁר לֹא דִבְּרוֹ יְהֹוָה׃}
}
{
	\eDtn{
		“And if you’ll say in your heart, ‘How shall we know the
		thing, that YHWH didn’t speak it?!’—
	}
}
{
%	\footnote{
%		How shall we know. It is one of the Bible’s central and most difficult
%		questions: How does one tell a true prophet from a false one? Moses
%		tells the people that the way to tell a false prophet is by seeing
%		whether his prophecy comes true or not. But that is a little late, is it
%		not? The question was how to know at the time of the prophecy whether it
%		is from God. Moses’ instruction appears to mean that one should go by
%		the prophet’s past record. Even then, people’s inclination seems to be
%		to disbelieve the true prophets. Even after Jeremiah’s prophecies of
%		Jerusalem’s fall come to pass (and the prophecies of those who oppose
%		him fail), as soon as he gives a prophecy that the people do not like
%		they say, “You’re speaking a lie. YHWH hasn’t sent you” (Jer 43:2)! So
%		Moses’ criterion for identifying false prophets may seem simple and
%		obvious, but the psychological point is that people miss the obvious and
%		turn instead to the comfortable.
%	}
}

\Verse{22}
{
	\hDtn{אֲשֶׁר יְדַבֵּר הַנָּבִיא בְּשֵׁם יְהֹוָה וְלֹא יִהְיֶה הַדָּבָר וְלֹא יָבֹא הוּא הַדָּבָר אֲשֶׁר לֹא דִבְּרוֹ יְהֹוָה בְּזָדוֹן דִּבְּרוֹ הַנָּבִיא לֹא תָגוּר מִמֶּנּוּ׃}
}
{
	\eDtn{
		when the prophet will speak in YHWH’s name, and the thing
		won’t be and won’t come to pass: that is the thing that
		YHWH did not say. The prophet spoke it presumptuously. You
		shall not be fearful of him.
	}
}
{
}
